% Source-faithful assembly: Mumford, Abelian Varieties, Chapter IV,
% Horizon transcription pages 0207--0236.

\section{Classification of positive involutions}
\label{mumford-IV21-classification-source}
\dcref{M:ch4:IV.21}

\begin{theorem}[Brauer groups and positive involutions]
\label{mumford-thm-positive-involution-classification}
For a division algebra $D$ with positive involution, the involution is one of
the Albert types: (I) a totally real field with identity involution; (II)
a quaternion division algebra over a totally real field with an orthogonal
involution; (III) a quaternion division algebra with its canonical
involution; or (IV) an algebra with totally imaginary quadratic center over
a totally real fixed field, with conjugation on the complex factors.
\source{mumford-abelian-varieties:page-0207}
\source{mumford-abelian-varieties:page-0208}
\source{mumford-abelian-varieties:page-0209}
\source{mumford-abelian-varieties:page-0210}
\source{mumford-abelian-varieties:page-0211}
\source{mumford-abelian-varieties:page-0212}
\source{mumford-abelian-varieties:page-0213}
\end{theorem}
\begin{proof}
For an involution of the second kind, class field theory gives
\[
0\to\operatorname{Br}(K)\to\prod_v\operatorname{Br}(K_v)
\to\mathbb Q/\mathbb Z\to0,
\]
with $\operatorname{Br}(\mathbb R)$ of order two and
$\operatorname{Br}(\mathbb C)=0$. If $\sigma$ is the involution on the
center, local invariants satisfy
$\operatorname{Inv}_v(D)+\operatorname{Inv}_{\sigma v}(D)=0$.
Skolem--Noether writes $x'=ax^*a^{-1}$, and the norm condition
$a^*a\in N_{K/K_0}K^*$ is equivalent to the existence of the involution.
For a first-kind involution, the Brauer class has order one or two; the
nontrivial case is quaternionic. Extending to the real factors, positivity
rules out $M_2(\mathbb R)$ in the canonical case and leaves Hamilton factors.
For complex factors the involution is $(X_i)\mapsto
A_i\overline{X_i}A_i^{-1}$; positivity makes $A_i$ positive definite,
and an inner change reduces to standard conjugation. These calculations give
the four listed types.
\end{proof}

\begin{proposition}[Albert numerical table]
\label{mumford-prop-albert-invariants}
For $S=\{x\in D:x'=x\}$, $\eta=\dim_{\mathbb Q}S/\dim_{\mathbb Q}D$,
and $[D:K]=d^2$, $[K_0:\mathbb Q]=e_0$, the four types have
\[
\begin{array}{c|c|c|c}
\text{type}&e&d&\eta\\ \hline
I&e_0&1&1\\ II&e_0&2&1/2\\ III&e_0&2&1/2\\ IV&2e_0&d&1/2.
\end{array}
\]
For endomorphism algebras of simple abelian varieties these invariants obey
the divisibility restrictions by $g=\dim X$.
\source{mumford-abelian-varieties:page-0212}
\source{mumford-abelian-varieties:page-0213}
% Source states the table and explicitly says the remaining restrictions are not proved here.
\notready
\end{proposition}

\section{The Riemann hypothesis}
\label{mumford-IV-rh-source}
\dcref{M:ch4:IV.21}

\begin{proposition}[Characteristic roots]
\label{mumford-prop-characteristic-root-modulus}
If $\alpha'\alpha=a\in\mathbb Z$ and $\omega_i$ are the roots of the
characteristic polynomial of $\alpha$, then $\mathbb Q[\alpha]$ is
semisimple, $|\omega_i|^2=a$, and $\omega_i\mapsto a/\omega_i$ permutes
the roots.
\source{mumford-abelian-varieties:page-0214}
\source{mumford-abelian-varieties:page-0215}
\end{proposition}
\begin{proof}
The minimal polynomial divides the characteristic polynomial and has the
same roots. The trace restricted to $\mathbb Q[\alpha]$ is positive; an
ideal and its trace-orthogonal complement split the algebra, proving
semisimplicity. Its simple factors are totally real fields or totally
imaginary quadratic extensions. The involution sends each root to its
complex conjugate, while $\alpha'\alpha=a$, hence
$a=\omega_i\overline{\omega_i}=|\omega_i|^2$.
\end{proof}

\begin{theorem}[Lang and Weil]
\label{mumford-thm-lang-frobenius-rh}
If $X_0$ is a scheme over $\mathbb F_q$ whose base change $X$ is an abelian
variety, then $X_0(\mathbb F_q)\ne\varnothing$. If $\pi$ is Frobenius and
$\omega_i$ are its characteristic roots, then
\[
\#X(\mathbb F_{q^n})=\prod_i(1-\omega_i^n),
\qquad |\omega_i|=\sqrt q,
\qquad \omega_{\pi(i)}=q/\omega_i.
\]
\source{mumford-abelian-varieties:page-0215}
\source{mumford-abelian-varieties:page-0216}
\source{mumford-abelian-varieties:page-0217}
\end{theorem}
\begin{proof}
Frobenius is the identity on the underlying space, is functorial, and has
zero differential. Write it as $\pi(x)=x_0+f(x)$. Then $1-f$ has identity
differential, hence zero-dimensional kernel and is surjective. Choose
$x_1$ with $(1-f)x_1=x_0$; then $x_1=\pi(x_1)$.
Thus $\#X(\mathbb F_{q^n})=\deg(1-\pi^n)$. After replacing the field,
choose an ample bundle; Rosati gives $\pi'\pi=q$, yielding the reciprocal
root relation and $|\omega_i|=\sqrt q$.
\end{proof}

\begin{corollary}[Weil estimate]
\label{mumford-cor-weil-point-count-estimate}
For some constant $C$,
$|\#X(\mathbb F_{q^n})-q^{ng}|\le Cq^{n(g-1/2)}$.
\source{mumford-abelian-varieties:page-0218}
\uses{mumford-thm-lang-frobenius-rh}
\end{corollary}

\begin{theorem}[Serre torsion restriction]
\label{mumford-thm-serre-torsion-restriction}
For $n\ge3$ and ample $L$, the restriction from automorphisms preserving
$L$ modulo $\operatorname{Pic}^0X$ to $\operatorname{Aut}(X_n)$ is injective.
\source{mumford-abelian-varieties:page-0218}
\end{theorem}
\begin{proof}
The characteristic roots are algebraic integers all of whose conjugates
have absolute value one, hence roots of unity. If $\alpha$ is the identity
on $X_n$, write $\alpha-1=n\beta$; the root-of-unity lemma forces
$\alpha=1$ for $n>3$.
\end{proof}

\section{The Jordan algebra of $NS^0(X)$}
\label{mumford-IV21-jordan-source}
\dcref{M:ch4:IV.21}

\begin{proposition}[Real structure of $NS^0$]
\label{mumford-prop-ns-jordan-structure}
For an ample $L$,
\[
NS^0(X)\simeq\{\alpha\in\operatorname{End}^0X:\alpha'=\alpha\}.
\]
With product $\alpha\circ\beta=(\alpha\beta+\beta\alpha)/2$, this is a
formally real Jordan algebra, a product over $\mathbb R$ of real symmetric,
complex Hermitian, and quaternionic Hermitian matrix algebras. Ample bundles
are exactly its totally positive elements.
\source{mumford-abelian-varieties:page-0219}
\source{mumford-abelian-varieties:page-0220}
\source{mumford-abelian-varieties:page-0221}
\uses{mumford-thm-ns-rosati-fixed,mumford-thm-rosati-positive}
% Source gives the application and numerical formulas but no proof in these pages.
\notready
\end{proposition}

\section{Complex multiplication and finite fields}
\label{mumford-IV22-source}
\dcref{M:ch4:IV.22}

\begin{theorem}[CM complex-torus classification]
\label{mumford-thm-cm-complex-classification}
Let $K_0$ be totally real of degree $g$ and $K$ a totally imaginary quadratic
extension. The pairs $(X,i)$ with $X/\mathbb C$ an abelian variety and
$i:K\hookrightarrow\operatorname{End}^0X$ have exactly $2^g$ isomorphism
classes up to isogeny. Representatives are
$X(K,\Phi)=(\mathbb R\otimes_{\mathbb Q}K)/A_0$, one for each of the
$2^g$ complex-algebra structures $\Phi$ on $\mathbb R\otimes K$.
\source{mumford-abelian-varieties:page-0221}
\source{mumford-abelian-varieties:page-0222}
\source{mumford-abelian-varieties:page-0223}
\source{mumford-abelian-varieties:page-0224}
\end{theorem}
\begin{proof}
The lattice $U$ is a finite-index $A_0$-module; after isogeny it is $A_0$.
The $K$-action pulls the complex structure back to a $\mathbb C$-algebra
structure $\Phi$. There are two embeddings over each real embedding of
$K_0$, hence $2^g$ structures. Choose $\alpha\in K$ with
$\alpha^2\in K_0$ and embeddings $\tau_i(\alpha)=i\beta_i$, $\beta_i>0$.
Then
\[
H(x,y)=2\sum_i\beta_i\tau_i(x)\overline{\tau_i(y)}
\]
is positive definite, and for $x,y\in A_0$,
\[
\operatorname{Im}H(x,y)=-\operatorname{Tr}_{K/\mathbb Q}(\alpha xy)\in\mathbb Z.
\]
Thus the torus is an abelian variety. An equivalence preserving the $K$-
action is $K$-linear and, after normalization, the identity; so distinct
$\Phi$ are inequivalent.
\end{proof}

\begin{theorem}[Deuring]
\label{mumford-thm-deuring-elliptic}
For an elliptic curve in characteristic $p$, non-finite-field definition is
equivalent to $\operatorname{End}^0X=\mathbb Q$. If $X$ is defined over a
finite field, its endomorphism algebra is imaginary quadratic exactly when
the $p$-rank is one; for $p$-rank zero it is the quaternion division algebra
with local invariants $\operatorname{Inv}_p=\operatorname{Inv}_\infty=1/2$
and all other finite invariants zero.
\source{mumford-abelian-varieties:page-0228}
\source{mumford-abelian-varieties:page-0229}
\source{mumford-abelian-varieties:page-0230}
\source{mumford-abelian-varieties:page-0231}
\end{theorem}
\begin{proof}
Let (A) be $p$-rank zero, (B) noncommutative endomorphism algebra, and (C)
finite-field definition with $\pi^n=p_X^m$. The classification table and
the $p$-adic Tate representation give (B)$\Rightarrow$(A). If (A) and the
algebra is commutative, cyclic Tate subgroups and quotients yield an
endomorphism with cyclic $l$-power kernel, contradicting the norm in a
number field; hence (A)$\Rightarrow$(B). Frobenius centrality gives
(A)$\Leftrightarrow$(C), and the local invariants follow from the sum of
invariants. In the non-finite-field case, a quadratic endomorphism acting
scalarly on a Tate module would have rational scalar with negative square,
impossible; hence the algebra is $\mathbb Q$.
\end{proof}

\begin{theorem}[CM over finite fields]
\label{mumford-thm-cm-finite-field-isogeny}
A simple abelian variety in characteristic $p$ is isogenous to one defined
over a finite field if and only if it is of CM type.
\source{mumford-abelian-varieties:page-0231}
% Source attributes the two implications to Tate and Grothendieck without reproducing proofs.
\notready
\end{theorem}

\begin{theorem}[Tate centralizer theorem]
\label{mumford-thm-tate-centralizer}
If $X$ is defined over a finite field $k_0$, with Frobenius $\pi$, then
\[
\operatorname{End}(X,k_0)\otimes\mathbb Q_l
=\operatorname{Cent}_{\operatorname{End}_{\mathbb Q_l}(T_lX)}(T_l\pi).
\]
\source{mumford-abelian-varieties:page-0232}
% Source introduces this as Tate's further theorem without a proof.
\notready
\end{theorem}

\section{The theta group of a line bundle}
\label{mumford-IV23-source}
\dcref{M:ch4:IV.23}

\begin{definition}[Theta group]
\label{mumford-def-theta-group}
A theta-group is an exact sequence
\[
1\to\mathbb G_m\xrightarrow{i}G\xrightarrow{\pi}K\to1
\]
where $K$ is commutative, the sequence is locally split, $i$ is a closed
immersion identifying the kernel, and $\mathbb G_m$ is central. For finite
$K$, a section identifies the underlying scheme with
$\mathbb G_m\times K$ and the group law with
\[
(\alpha,k)(\alpha',k')=(\alpha\alpha'f(k,k'),k+k'),
\]
where $f$ is a 2-cocycle, defined up to coboundary.
\source{mumford-abelian-varieties:page-0232}
\source{mumford-abelian-varieties:page-0233}
\end{definition}

\begin{proposition}[Theta commutator]
\label{mumford-prop-theta-commutator}
The commutator defines a skew-symmetric bihomomorphism
$e:K\times K\to\mathbb G_m$, and hence a homomorphism
$\gamma:K\to\widehat K$ with
$e(k,k')=\langle k,\gamma(k')\rangle$. The center of $G$ maps to
$\ker\gamma$; if $K$ is finite, the center is exactly $\mathbb G_m$ iff
$\gamma$ is an isomorphism.
\source{mumford-abelian-varieties:page-0233}
\source{mumford-abelian-varieties:page-0234}
\end{proposition}
\begin{proof}
The commutator is independent of lifts because $\mathbb G_m$ is central.
The group identities give bihomomorphy and skew-symmetry. Viewing
$(e,p_2):K\times K\to\mathbb G_m\times K$ over $K$ gives $\gamma$.
The displayed universal pairing shows that $\gamma(\pi(x))$ is trivial
exactly when $x$ commutes with every base change of every point of $G$.
\end{proof}

\begin{lemma}[Trivial theta extensions]
\label{mumford-lem-theta-extension-prime-order}
If $K$ is finite and $G$ commutative, then
$G\simeq\mathbb G_m\times K$ as a group. If $K$ has prime order, then $G$
is commutative.
\source{mumford-abelian-varieties:page-0235}
\end{lemma}
\begin{proof}
The first assertion reduces by an Ext sequence to cyclic prime-order,
$\mu_p$, and $\alpha_p$ cases; divisibility of $k^*$ splits the reduced
cases and the height-one Lie algebra calculation splits the latter cases.
For prime order, reduced cases are generated by a central subgroup and one
lift. In the infinitesimal cases the Lie algebra is abelian, so the
height-one subgroup is commutative; together with the central torus this
makes $G$ commutative.
\end{proof}

\begin{theorem}[Theta group of a line bundle]
\label{mumford-thm-line-bundle-theta-group}
For a line bundle $L$ on an abelian variety $X$, the functor of automorphisms
of $S\times L$ covering translations of $S\times X$ is represented by a
group scheme $\mathcal G(L)$, fitting into a theta-group extension
\[
1\to\mathbb G_m\to\mathcal G(L)\to K(L)\to1.
\]
\source{mumford-abelian-varieties:page-0235}
\source{mumford-abelian-varieties:page-0236}
\uses{mumford-def-theta-group}
\end{theorem}
\begin{proof}
Let $L^*$ be the complement of the zero section. Fix $P_0$ over the origin
and set $\mathcal G(L)=p^{-1}(K(L))\cap L^*$. An automorphism over a
translation $T_f$ sends $P_0$ to a morphism $S\to L^*$ over $f$, hence
factors through $\mathcal G(L)$. This gives the representing map; its
inverse is obtained by translating the fibre and using linearity. Scalar
automorphisms form $\mathbb G_m$, and the quotient is $K(L)$.
\end{proof}
